\chapter{Bed Bugs}
\index{Bed Bugs}

\section*{Definition and Overview}
Bed bugs are small, wingless insects that feed on the blood of humans and animals, usually at night while people sleep. They are not known to transmit infectious diseases, but their bites can cause itchy skin reactions, and an infestation can cause significant distress, sleep disturbance, and anxiety. Bed bugs are not a sign of poor hygiene; they can be picked up in any home, hotel, or form of public transport and are increasingly common worldwide.

\section*{Epidemiology}
Bed bug infestations have increased markedly in many countries over recent decades, linked to increased international travel, resistance to insecticides, and changes in pest-control practices. Bed bugs affect people of all ages and both sexes, in all socioeconomic groups and types of housing. They are most common where there is high occupancy turnover, such as hotels, hostels, apartments, and laundromat or second-hand furniture sharing. Outbreaks are common in student accommodation and aged-care settings.

\section*{Aetiology and Pathophysiology}
\begin{itemize}
    \item \textbf{The insect:} the common bed bug (\textit{Cimex lectularius}) feeds on blood using a sharp proboscis that pierces the skin
    \item \textbf{Feeding:} bed bugs usually feed at night, are attracted by body heat and carbon dioxide, and can hide for many months between feeds
    \item \textbf{Bite reaction:} the insect injects saliva containing anticoagulants and other substances while feeding; the body's immune response to this saliva causes the itchy bumps
    \item \textbf{Spread:} bed bugs travel by hitch-hiking in luggage, clothing, bedding, furniture, and other belongings; they do not fly and do not live on people
    \item \textbf{Hiding places:} cracks in beds, mattress seams, skirting boards, electrical fittings, and any crevice near the bed
\end{itemize}

\section*{Clinical Features}
\begin{itemize}
    \item Small, red, itchy bumps on the skin, often in a line or cluster (sometimes described as breakfast, lunch, and dinner)
    \item Bites most commonly on exposed skin: face, neck, arms, and hands
    \item Itching may be intense, and some people develop blistering
    \item Not everyone reacts; some people show no marks at all, while others develop delayed reactions days later
    \item Blood spots on sheets, dark rust-coloured specks of droppings on bedding and walls, and a sweet, musty odour in heavily infested rooms
\end{itemize}

\section*{Differential Diagnosis}
\begin{itemize}
    \item \textbf{Flea bites:} typically around the ankles and lower legs
    \item \textbf{Mosquito bites:} single, scattered, itchy bumps
    \item \textbf{Mite bites and scabies:} intense itching, often between fingers and in skin folds, with burrows
    \item \textbf{Other insect bites,} such as spiders or lice
    \item \textbf{Allergic skin conditions,} such as urticaria (hives) or eczema
\end{itemize}

\section*{When to Seek Medical Care}
See a doctor if bite marks are very itchy, become painful, or show signs of infection such as spreading redness, swelling, or discharge. Most people manage bite symptoms at home.

\textbf{Contact your local emergency services for:}
\begin{itemize}
    \item Signs of a severe allergic reaction, such as difficulty breathing, swelling of the face or lips, or dizziness
    \item Rapidly spreading redness, pain, or fever suggesting skin infection
    \item Any severe reaction that worries you
\end{itemize}

\section*{Diagnosis and Investigations}
\begin{itemize}
    \item \textbf{History:} itchiness, sleep disturbance, recent travel, second-hand furniture, or known infestation at home, work, or school
    \item \textbf{Examination:} inspection of the bites and their pattern (lines or clusters on exposed skin)
    \item \textbf{Search for the insects:} the diagnosis is confirmed by finding live bed bugs, shed skins, or droppings in the mattress seams, bed frame, and nearby crevices
    \item \textbf{No blood tests:} there is no useful blood test for bed bug bites
    \item \textbf{Pest-control assessment:} a professional may be needed to confirm and locate the infestation
\end{itemize}

\section*{Management}
\begin{itemize}
    \item \textbf{Bite care:} wash bites gently with soap and water; apply a cool compress to reduce itching
    \item \textbf{Anti-itch measures:} calamine lotion, an over-the-counter corticosteroid cream, or an oral antihistamine can relieve itching
    \item \textbf{Avoid scratching:} keep nails short to reduce skin damage
    \item \textbf{Treatment of infection:} antibiotics if a bite becomes infected
    \item \textbf{Pest eradication:} professional pest control using approved insecticides, plus heat treatment (steam or laundering at high temperature) of bedding and clothing
    \item \textbf{Follow-up:} repeated inspections and treatments may be needed because bed bug eggs are resistant to many chemicals
\end{itemize}

\section*{Complications}
\begin{itemize}
    \item Secondary bacterial infection of scratched bites, such as impetigo or cellulitis
    \item Scarring or persistent skin discoloration after healing
    \item Sleep deprivation, anxiety, and psychological distress
    \item Rarely, severe allergic reactions to bites
    \item There is no evidence that bed bugs transmit infectious diseases to humans
\end{itemize}

\section*{Prognosis and Outcomes}
Bed bug bites generally heal within one to two weeks without any treatment, although itching can last longer. Infestations are treatable, but eradication may take several weeks and require persistence and repeated professional treatment. The insects do not live on the body, so personal hygiene does not need to change. Once the home is cleared, there is usually no lasting health effect beyond temporary skin and psychological disturbance.

\section*{Prevention and Self-Care}
\begin{itemize}
    \item Inspect mattresses, bed frames, and second-hand furniture for signs of bed bugs before bringing items home
    \item When travelling, check hotel bedding and keep luggage off the floor and on a rack
    \item After travel, inspect and vacuum luggage and wash clothes on a hot cycle
    \item Use protective mattress and pillow encasements
    \item Reduce clutter in the bedroom to remove hiding places
    \item If infested, launder bedding and clothing at 60 degrees or above, vacuum thoroughly, and arrange professional pest control
    \item Do not move infested furniture into other rooms or discard it where others may take it, as this spreads the problem
\end{itemize}

\section*{Sources and Further Reading}
The following sources provide reliable, up-to-date information on bed bugs and bite management.
\begin{itemize}
    \item Centers for Disease Control and Prevention. \textit{Bed Bugs.} https://www.cdc.gov
    \item NHS. \textit{Bedbugs.} https://www.nhs.uk/conditions/
    \item Mayo Clinic. \textit{Bedbugs.} https://www.mayoclinic.org
    \item MedlinePlus. \textit{Insect Bites and Stings.} https://medlineplus.gov
    \item World Health Organization. \textit{Vector-borne diseases.} https://www.who.int
\end{itemize}
